German election files contain two separate ballots per election year:
\textit{Erststimme} (direct constituency vote) and
\textit{Zweitstimme} (party list vote). The \emph{Zweitstimme} determines
the proportional allocation of seats in the Bundestag, while the \emph{Erststimme} 
determines which candidate wins the direct constituency mandate.  

Results are reported at the \textit{Gemeinde} level where available, with
postal votes (\textit{Briefwahl}) included as special units attached to their
respective district (\textit{Landkreis}, NUTS~3). For the 2004 EP election only,
results are reported at the \emph{Landkreis} level, since municipality-level 
data is not publicly available.
